\documentclass[../main.tex]{subfiles}
\graphicspath{{\subfix{../images/}}}
\begin{document}

\definecolor{hwrhell}{RGB}{248,232,232}
\definecolor{linie}{RGB}{120,120,120}

Die ausgewählte Lösung (Gesamtlösung 1, Kapitel 7) ist im Folgenden als Zeichnungssatz dokumentiert. Grundlage ist das 3D-Modell Mikrocontrollergehäuse.step. Alle Maße wurden direkt aus der Geometrie des Modells ermittelt. Der Satz besteht aus der Übersichts- bzw. Explosionsansicht des Gehäuses inklusive Stückliste (Tabelle \ref{tab:stueckliste}) und den Einzelteilzeichnungen der beiden Gehäuseteile (Gehäuseoberteil und Gehäuseunterteil).

\begin{table}[htbp]
\centering
\caption{Stückliste SCARA-DJ01-Gehäuse}
\label{tab:stueckliste}
\small
\renewcommand{\arraystretch}{1.25}
\begin{tabularx}{\textwidth}{@{}c c >{\raggedright\arraybackslash}p{4.1cm} >{\raggedright\arraybackslash}p{3.3cm} >{\raggedright\arraybackslash}X@{}}
\toprule
\rowcolor{hwrhell}
\textbf{Pos.} & \textbf{Menge} & \textbf{Benennung} & \textbf{Werkstoff / Art.-Nr.} & \textbf{Abmessung / Bemerkung} \\
\midrule
1 & 1 & Leiterplatte, bestückt & FR-4, 1,6\,mm & $90\times83\times1{,}6$ \\
2 & 6 & Senkkopfschraube M3 & Edelstahl AISI 304/ Art.-Nr. 92010A126 & Gesamtlänge 16\,mm; verbindet Gehäuseunter- und Gehäuseoberteil \\
3 & 10 & Gewindeeinsatz M3 (Einschmelzgewinde) & Messing & Außen-Ø 5\,mm, Gesamtlänge 5\,mm; 6$\times$ in Oberschale, 4$\times$ in Unterschale\\
4 & 6 & Zylinderschraube mit niedrigem Kopf M3 & Edelstahl AISI 304/ Art.-Nr. 92855A307 & Gesamtlänge 6\,mm; verbindet PCB und Gehäuseunter \\
5 & 1 & Gehäuseunterteil & Kunststoff/PLA (FDM) & $100\times253\times42$ \\
6 & 1 & Gehäuseoberteil & Kunststoff/PLA (FDM) & $100\times249\times67$ \\
\bottomrule
\end{tabularx}
\end{table}

\clearpage
\newgeometry{a2paper, margin=1.5cm}
\begin{landscape}
  \centering
  \includegraphics[width=\linewidth, height=0.9\textheight, keepaspectratio]{Bilder/TZ-S1.png}
  \captionof{figure}{Technische Zeichnung Blatt 1/3: Übersichts-/Explosionsansicht mit Stückliste}

  \clearpage
  \centering
  \includegraphics[width=\linewidth, height=0.9\textheight, keepaspectratio]{Bilder/TZ-S2.png}
  \captionof{figure}{Technische Zeichnung Blatt 2/3: Gehäuseunterteil}

  \clearpage
  \centering
  \includegraphics[width=\linewidth, height=0.9\textheight, keepaspectratio]{Bilder/TZ-S3.png}
  \captionof{figure}{Technische Zeichnung Blatt 3/3: Gehäuseoberteil}
\end{landscape}
\restoregeometry

\end{document}